\begin{table}[htbp]
\centering\small
\begin{tabularx}{\linewidth}{rlXrrrrr}
\toprule
順 & 型 & テーマ & 問いの価値 & 信頼性 & 実行可能性 & 競争と出口 & 合計 \\
\midrule
1 & 手法 & LLM が作った変数を回帰に使うときの測定誤差の補正 & 70 & 87 & 85 & 70 & 77.5 \\
2 & 手法 & 仕様曲線と頑健性分析の自動化、多重性を考えた推論 & 60 & 80 & 85 & 70 & 73.1 \\
3 & 因果 & 処置時点がずれる差の差研究の大規模な再推定 & 60 & 80 & 85 & 65 & 71.8 \\
4 & 手法 & 処置時点がずれる差の差推定量の統一的な実装と診断 & 55 & 93 & 80 & 55 & 68.9 \\
5 & 構造 & 動学的離散選択モデルの推定を速くする近似手法 & 55 & 80 & 75 & 65 & 68.1 \\
6 & 手法 & クラスタ数が少ないときの推論手法の高速実装と比較 & 50 & 87 & 80 & 55 & 66.1 \\
7 & 構造 & 製品の文章と画像から特性を作る需要推定 & 65 & 60 & 65 & 75 & 66.0 \\
8 & 因果 & 生成 AI の普及が職種別の求人と賃金に与える効果 & 65 & 53 & 70 & 60 & 61.8 \\
9 & 因果 & 地域別最低賃金の引き上げが雇用と賃金分布に与える効果 & 75 & 67 & 60 & 40 & 58.9 \\
10 & 構造 & 公開データによるオンライン市場の価格設定の構造推定 & 60 & 60 & 60 & 55 & 58.7 \\
11 & 因果 & 在宅勤務の普及が賃金と地域間移動に与えた効果 & 60 & 53 & 55 & 45 & 53.0 \\
12 & 構造 & 職探しモデルによる失業給付の設計 & 70 & 60 & 45 & 40 & 52.4 \\
\bottomrule
\end{tabularx}
\caption{星取り表（合計点の順）。群の点は20〜100、合計は4群の点の幾何平均。12件中12件が試行採点（出典未調査）、11件は根拠の欄が空で、新規性かデータ入手性を上限2点で計算した。データは \texttt{theme/research/themes.csv}。}
\label{tab:th-hoshi}
\end{table}
